\documentclass[12pt]{article}

\usepackage[a4paper, total={6in, 8in}]{geometry}
\usepackage{textcomp}

\begin{document}
\setlength{\parindent}{0pt}

\def \t {\textrightarrow}
\def \v {\vspace{1em}}
\def \bi {\begin{itemize}}
\def \ei {\end{itemize}}
\def \s[#1] {\section*{#1}}
\def \ss[#1] {\subsection*{#1}}
\def \sss[#1] {\subsubsection*{#1}}

\s[Bergson - Le due fonti della morale e della religione]
Bergson fa anche un'altra riflessione oltre al discorso metafisico, scrive "Le due fonti della morale e della religione" del 32 \t è la sua ultima opera \\
Diventa famosa perchè si interessa dal punto di vista etico-politico e religioso della struttura sociale \t è opera filosofica ma anche di scienza sociale \\
Lo slancio vitale che nelle altre speci animali si ferma, nell'uomo è ancora presente \t e si esprime in tutta l'attività creativa dell'uomo \t in tutte le forme che sono espressione dell'uomo (filosofia, arte, religione) \\
Che sono realta non di natura, ma dipendono dal fatto che l'uomo è in continua evoluzione \\
Si sofferma in particolare sulla creativià morale e religiosa \t sposta l'oggetto del suo interesse quindi \\
Parte dallo studio della cosceinza, con "evoluzione creatrice" passa a teoria dell'universo, mentre in quest'opera elabora teoria dei valori

\ss[Morale]
\ss[Pressione sociale]
Ci sono delle norme morali che nascono dal fatto che vita in comunità ha delle esigenze \t ci sono delle regole che uomo si è dato \\
Per stare in quel gruppo io devo sottostare a questa norma \t e questo lo insegna la storia \t ogni uomo del suo tempo è stato nella società, e ha fatto proprie le norme sociali \\
In questo gruppo sociale ripercorro quindi la strada di altri \t non invento niente di nuovo, mi adeguo solo alla vita del gruppo \t e mi abituo quindi \\
Mi abituo a contrarre abitudini \t mi abituo a cio che in quel gruppo sociale viene considerato valore \\
La morale che nasce a opera della pressione sociale si fonda sull'obbligo di stare in quelle regole, se voglio stare in quel gruppo sociale \\
No atteggiamento di disprezzo, ma rilevazione di un fatto da parte di Bergson \t se voglio stare in quel gruppo, devo rispettare quelle regole \t se no mino la stabilità del gruppo \\
La morale è della societa chiusa \t io agisco come parte del tutto, e il tutto puo essere:
\bi
  \item la nazione
  \item una associazione
  \item la famiglia
\ei
Mi conformo alle regole relative al determinato gruppo \\
Questo tipo di morale assicura la vita della comunità \t il rispetto delle regole permette la sopravvivenza del gruppo

\ss[Slancio d'amore]
Esite anche una morale aperta pero, una morale della società aperta, assoluta \t che rompe il gruppo, rompe l'abitudine \t è una morale creatrice \\
È una morale che spezza l'equilibrio del gruppo, e possono essere:
\bi
\item del cristianesimo
\item saggi della Grecia
\item profeti di Israele
\ei
Socrate come gesu ha creato valori totlamente diversi dal gruppo di appartenenza \t hanno messo in discussione il gruppo aggiungendo un valore aggiunto \\
La societa aperta è l'umanita \t i valori della societa aperta non si rivolgono a un gruppo, ma all'umo quanto tale \\
Quindi la morale della società aperta è dinamica \t non è come quella chiusa, che mette dei valori ai quali bisogna omologarsi e rimangono stabili \t non c'è creazione ne adeguamento nuovo, solo abitudine \\
Finora solo società chiuse \t pero bisogna tendere verso società aperta \t bisogna realizzare una morale aperta, definire dei valori universali che valgano oltre il limite del gruppo \\
Bisogna rompere il limite dell'associazione \\
C'è però un salto qualitativo per passare a societa aperta \t non è un passaggio ?? \\
Il fondamento di questa morale aperta è la persona, creatrice di valori \t e il fine è l'umanità intera, e mi rivolgo a tutti gli uomini in modo innovativo, rompendo tutti gli schemi delle società chiuse \\
Il contenuto fondamentale della morale deve essere l'amore per l'uomo \t la morale aperta deve considerare l'uomo in quanto tale \\
Questa morale però non si puo insegnare \t Socrate, Gesu non l'hanno imparata a scuola \t è la morale di coloro che rompono gli schemi, che seguono la loro intuizione \\
Il limite di questa morale è che non si può insegnare

\ss[Religione]
\ss[Religione statica]
= sito della funzione fabulatrice \t religione intessuta di miti e di favole, di racconti \\
È la forma di religiostia che si è formata durante l'evoluzione per scopi principalmente di soppravvivenza \t l'uomo è intelligente, e fa paura \t puo pensare qualsiasi cosa e puo volgere le sue idee anche vero la non vita \\
Quindi bisogna contenere l'impredivibilita dell uomo intelligente \t e l'uomo è anche fondamentalmente egoista \t ci sono dei pericoli che caratterizzano umano \\
Quindi storicamente questo tipo di religione è stato utile per contenere questa possibile pericolostia dell uomo e rafforzare relazioni sociali \\
Inoltre da speranza dell'immortalita e di protezione sovrannaturale \t per rendere l'uomo piu docile, per farlo vivere in una società \t e per evitare che l'uomo potesse cominciare a pensare solo a se stesso, a pensarsi come un individuo al di fuori di un gruppo \\
La religione statica è una difesa contro la minaccia dell'intelligenza umana \\
È pero frutto anche di una evoluzione naturale \t è una religione riduttiva rispetto a capacità razionale dell'uomo \t sta in mezzo alle intelligenze, infraintellettuale

\ss[Religione dinamica]
Va oltre i dogmi, che sono delle cristallizzazioni della religione stessa \t e lo immette nello slancio vitale \\
La religione dinamica è il misticismo \t i grandi mistici l'hanno vissuta, misticismo = presa di contatto tra uomo e Dio \\
Non solo l'uomo va a Dio, ma anche il contrario \t anche Dio fa sforzo, e questo rompe gli schemi \t il mistico non è interprete di una religione, ma incontra Dio direttamente \t senza religione \\
Es. San Francesco D'assisi, Giovanna d'arco, etc. \t il contatto superiore con Dio è il motivo per cui sono tornati nel mondo \t non è un misticismo dell'eremita quindi, ma del riportare nel mondo e nel quotidiano Dio \\
L'amoore per dio del mistico diventa quindi amore per l'umanità \t e il misticismo è l'unica prova dell'esistenza di Dio \\
Trova tra i mistici di tutte le religioni dei punti di tangenza, quindi dio esiste \\
L'umanita oggi pero non ha mistici, ma ne avrebbe grande bisogno


\s[Benedetto Croce]
Napoli è la culla dell'idealismo italiano \t e qua c'è Bertrando Spaventa è stato uno dei grandi protagonisti della diffusione dell hegelismo in Italia \\
Cerchera anche di dare un taglio + teoretico dell hegelismo \\
Ha molti seguaci, tra cui Francesco De Santis (1800) \t è un grande critico letterario \t scrive storia della letteratura italiana (1872), ispirandosi al concetto hegeliano di spirito \\
Lo spirito attraverso la coscienza prende conoscenza di se \t la poesia è manifestazione dello spirito \\
De Santis scrive anche "Leopardi e Schopenhauer" 

\v

Croce legge Spaventa, ma questa lettura lo allontana da Hegel \t ci arriva invece attraverso Marx \\
Croce nasce ad Aquila nel 1866 in una famiglia ricca e conservatrice \t studia a Napoli \\
Nel 83 si trova in vacanza a Ischia, e un terremoto distrugge l'isola e uccide tutta la sua famiglia \t Croce rimane anche lui sepolto per ore \t e viene affidato allo zio Silvio Spaventa, fratello di Bertrando Spaventa
\end{document}
